\input{./._relpath-to-latexroot.ltx}
\documentclass[\latexroot/\projectname]{subfiles}
\input{\latexroot/@resources/subfile-setup.ltx}

\begin{document}

\section{Introduction}
\whenintegrated{\label{sec:intro}} % integrated doc: SST for labels
\setcounter{page}{0}\pagenumbering{arabic}

High-frequency changes in interest rates around central bank announcements—commonly referred to as monetary policy or interest rate surprises—have become a widely used proxy for monetary policy (MP) shocks. A growing literature combines these surprises with vector autoregressive (VAR) models or local projections to analyze the macroeconomic effects of monetary policy, primarily in advanced economies (AEs). Applying this methodology to emerging market economies (EMEs), however, poses additional challenges, partly reflecting the more limited availability and shorter span of high-frequency financial data. 

Although recent contributions have extended high-frequency identification (HFI) techniques to some EMEs, a number of studies document puzzling responses to contractionary MP shocks, including exchange-rate depreciation, rising prices, and increases in inflation expectations. These findings are difficult to reconcile with standard macroeconomic models and raise important questions about the transmission and identification of MP shocks in EMEs.

A potential explanation for these puzzling findings is that interest rate surprises may not exclusively capture exogenous deviations from a systematic policy rule. If market participants are imperfectly informed about the state of the economy, these surprises may be contaminated by \emph{central bank information} (CBI) effects (\hyperlink{Romer and Romer (2000)}{Romer and Romer, 2000}; \hyperlink{Jarociński and Karadi (2020)}{Jarociński and Karadi, 2020}). Central banks may convey new information about the economic outlook in their announcements while simultaneously adjusting the policy stance in response. The resulting interest rate surprises then partly reflect an endogenous policy response to economic conditions, potentially biasing the transimission of MP shocks.

An additional source of contamination arises from \emph{market misperceptions about the central bank’s reaction function}. In this case, interest rate surprises may reflect an unexpectedly aggressive response of policymakers to observable developments, or \emph{central bank response-to-news} (RN) effects, as emphasized by \hyperlink{Bauer and Swanson (2023a, 2023b)}{Bauer and Swanson (2023a, 2023b)}. As shown by these authors, failing to control for these effects may also bias the transmission of MP shocks.


Recent work for the U.S. by \hyperlink{Jarociński and Karadi (2025)}{Jarociński and Karadi (2025)} jointly analyzes the role of both contamination channels. Their results indicate that CBI effects substantially attenuate the transmission of monetary policy, while RN effects appear to play a more limited role. We extend this analysis to an EME and jointly examine how CBI and RN components shape the transmission of two types of interest rate surprises. Specifically, we distinguish between surprises to the target policy rate, representing conventional MP shocks, and surprises to the expected near-term policy path, capturing unconventional or forward-guidance (FG) shocks.


Our analysis shows that both CBI and RN components appear to play an important role in attenuating the effects of target-rate surprises and in accounting for several of the puzzles documented in the EME literature. For FG surprises, however, RN components appear to play a more prominent role. Notably, RN contamination seems to operate asymmetrically across policy instruments: it weakens the transmission of target-rate shocks but disproportionately amplifies and accelerates the estimated effects of FG innovations.

Focusing on Mexico as an ideal case study, we begin by computing high-frequency changes in interest rate swaps at different maturities around monetary policy announcements.\footnote{Mexico provides an ideal case study given its inflation-targeting regime, a flexible and market-determined exchange rate, and relatively deep and liquid financial markets, which allow the measurement of monetary policy surprises across different maturities—an essential feature for distinguishing between conventional and unconventional MP shocks.} Following the standard approach of \hyperlink{Gürkaynak et al. (2005)}{Gürkaynak et al. (2005)}, we summarize these variations into a target-rate surprise, capturing unexpected changes in the policy rate, and an FG surprise, reflecting revisions to the expected near-term policy path.

Following the logic of \hyperlink{Bauer and Swanson (2023a, 2023b)}{Bauer and Swanson (2023a, 2023b)}, we then regress each policy surprise on a set of macroeconomic and financial variables observable to market participants prior to the announcement. The fitted value captures the portion of the surprise that may reflect an unexpectedly aggressive response of the central bank to observable economic developments. The residual isolates the unpredictable part of the surprise and therefore provides a cleaner proxy for exogenous monetary policy innovations.

In a second step, we exploit the high-frequency response of stock prices around policy announcements to further refine this decomposition and distinguish between the three shocks of interest for each surprise. We then trace their dynamic effects using a monthly Bayesian VAR estimated over the period January 2002 to December 2024.

In particular, a \emph{pure MP shock} is characterized by a negative co-movement between the residual component of the interest rate surprise and the stock price surprise. This pattern is consistent with an unexpected policy tightening that raises discount rates and lowers expected future cash flows, leading to an immediate decline in equity prices.

An \emph{RN shock}, in turn, is characterized by a negative co-movement between the fitted component of the interest rate surprise and the stock price surprise. This configuration captures episodes in which the central bank responds more aggressively than markets anticipated to publicly available information. Because that information was largely incorporated into expectations prior to the announcement, the only new element revealed within the announcement window is a stronger-than-expected policy response, which depresses equity prices through standard discount-rate and cash-flow channels.

Finally, a \emph{CBI shock} is associated with a positive co-movement between the interest rate surprise and the stock price surprise. This pattern may arise when the central bank reveals favorable new information about economic fundamentals while simultaneously tightening policy. In such cases, equity prices increase in response to improved growth prospects and higher expected future cash flows, despite the policy tightening.

Our main results indicate that CBI and RN components tend to attenuate—and even reverse—the effects of target-rate surprises. When these components are not disentangled, impulse responses to contractionary target-rate shocks exhibit the well-known output, price, and exchange-rate puzzles, with inflation expectations rising on impact. Once the CBI and RN components are isolated, however, these puzzles largely disappear: output and prices display a U-shaped response with a delayed trough, the domestic currency appreciates against the U.S.\ dollar, and inflation expectations remain well anchored. In addition, stock prices decline, and both the one-year and the ten-year government bond yields increase, with the long-term yield rising by a smaller amount.

In line with these results, favorable CBI shocks identified from target-rate surprises are followed by increases in interest rates, an expansion in output, higher prices, and rising inflation expectations. Notably, the Mexican peso depreciates against the U.S.\ dollar in response to these shocks, a pattern that may be consistent with upward revisions in inflation expectations, as emphasized by \hyperlink{Gurkaynak et al. (2021)}{G\"urkaynak et al. (2021)} and \hyperlink{Stavrakeva and Tang (2019)}{Stavrakeva and Tang (2019)}. These responses are consistent with an interpretation in which the central bank reveals favorable information about the economic outlook while simultaneously tightening policy to lean against emerging inflationary pressures, potentially operating through aggregate demand effects and exchange-rate pass-through.

In turn, hawkish RN shocks identified from target-rate surprises are followed by increases in interest rates, an expansion in output, higher prices, rising inflation expectations, and a depreciation of the domestic currency. These dynamics closely resemble those observed following favorable CBI shocks and are in line with \hyperlink{Bauer and Swanson (2023a, 2023b)}{Bauer and Swanson (2023a, 2023b)}. In general, they are consistent with a situation in which the central bank tightens policy more aggressively than markets anticipated in response to improving economic conditions that are publicly observable. In this case, the contractionary effects of tighter monetary policy are more than offset by the expansionary forces associated with the underlying macroeconomic developments.

Turning to the analysis of unconventional MP shocks, our results indicate that CBI effects appear to play a limited role in contaminating the transmission of FG surprises. In contrast, RN shocks tend to distort the estimated effects of these surprises by amplifying and accelerating their impact. When RN shocks are not disentangled, FG innovations appear to generate disproportionately large effects on macroeconomic variables, a pattern that may be consistent with what has been termed the \emph{forward guidance puzzle} in the literature (see \hyperlink{Del Negro et al. (2023)}{Del Negro et al., 2023}). Output, for example, declines significantly on impact. In addition, the magnitude of the contraction one month after the shock is roughly 50\% larger than the peak decline estimated following a pure target-rate shock, which occurs only after about 15 months. These results are difficult to reconcile with the conventional view that monetary policy affects real activity with substantial lags. They also contrast with the empirical evidence for the U.S. reported by \hyperlink{Swanson (2021, 2024)}{Swanson (2021, 2024)}, who finds that FG shocks generate macroeconomic effects that are comparable to—or even smaller than—those of conventional MP shocks.

Once RN shocks are isolated, however, the effects of FG innovations become more conservative and broadly similar to those of pure target-rate shocks. One plausible interpretation is that, following episodes in which economic activity fails to cool as expected, the central bank may adopt a more hawkish forward guidance stance to support the convergence of inflation toward target. In such circumstances, hawkish RN shocks may capture not only a tightening implied by the policy rule, but also the economy’s ongoing contractionary adjustment already underway. In a VAR that treats raw FG surprises as exogenous shocks, this combination can generate large and immediate contractionary responses, thereby overstating the effects of FG.

The remainder of the paper is organized as follows. \hyperlink{Section 2}{Section~2} reviews the related literature. \hyperlink{Section 3}{Section~3} outlines the construction of monetary policy surprises and the identification strategy used to isolate the shocks of interest. \hyperlink{Section 4}{Section~4} presents the VAR and the data used in the estimation. The empirical results are reported in \hyperlink{Section 5}{Section~5}. \hyperlink{Section 6}{Section~6} concludes.

% Smart bibliography: Only include bibliography if standalone AND has citations
\smartbib

\end{document}
